\documentclass[12pt]{article}
\usepackage[a4paper,margin=1in]{geometry}
\usepackage[T1]{fontenc}
\usepackage{lmodern,microtype}
\usepackage{amsmath,amssymb,amsthm,mathtools}
\usepackage{booktabs,array,enumitem,xcolor}
\usepackage[numbers,sort&compress]{natbib}
\usepackage[colorlinks=true,linkcolor=blue!55!black,citecolor=blue!55!black,urlcolor=blue!55!black]{hyperref}
\linespread{1.07}

\newtheorem{proposition}{Proposition}[section]
\theoremstyle{remark}
\newtheorem{remark}[proposition]{Remark}
\newcommand{\norm}[1]{\left\lVert#1\right\rVert}

\title{Physical Residues and Spectral Transversality\\
\large Conditioning of the Canonical Edge Quartet}
\author{Liang Wang}
\date{July 2026}

\begin{document}
\maketitle

\begin{abstract}
RH-194 identifies one physical edge quartet on each side of the
$\sigma=0.01$ model, and RH-196 proves that the corresponding
source--observation Riesz channels admit an exact optimally balanced packet.
This paper measures the physical geometry of those packets.

All eight unique modes have nonzero transfer residue; the minimum residue
modulus is $1.107\times10^{-2}$.  The canonical minimum cross singular value
is $5.207\times10^{-3}$ on the left and $1.052\times10^{-3}$ on the right.
Therefore the optimal biorthogonal frame norm products are respectively
$192.05$ and $950.26$.  These are finite transverse packets, but they are
not well conditioned.

The large condition is not merely a defective temporal coordinate choice.
At the latest accepted windows, the temporal oblique condition numbers are
within two percent of the canonical spectral optima.  Simultaneously the
temporal-to-spectral principal-angle gaps are below $0.051$ on all four
right/left comparisons.  Thus the late clock is converging to both the
physical subspaces and their intrinsic oblique geometry.

This is a finite positive result with a sharp warning: no all-level residue
or cross-angle lower bound is known.  Any later determinant route should
seek gauge-invariant cancellation of the large condition rather than claim
uniformly benign coordinates.
\end{abstract}

\section{Data and quantities}

The input is the frozen RH-194 eigendecomposition and root matching
\citep{WangRH194}.  For each side there are four simple physical modes with
source and observation states
\begin{equation}\label{eq:states}
 X_i=P_iS,
 \qquad Y_i=P_i^*O^*,
\end{equation}
and transfer residues
\begin{equation}\label{eq:residues}
 r_i=\langle Y_i,X_i\rangle_F.
\end{equation}

Let $Q_R,Q_L$ be orthonormal bases of the two four-dimensional spans.  The
cross singular values are those of
\begin{equation}\label{eq:cross}
 H=Q_L^*Q_R.
\end{equation}
Write
\begin{equation}\label{eq:gamma}
 \gamma=\sigma_{\min}(H).
\end{equation}
By RH-196, the optimal product of norms of biorthogonal frames on these
spaces is exactly
\begin{equation}\label{eq:condition}
 \chi_{\rm can}=\gamma^{-1},
\end{equation}
and each balanced frame has norm $\gamma^{-1/2}$
\citep{WangRH196}.

\section{Residue audit}

The unique-mode residue ranges are:
\begin{center}
\begin{tabular}{lrrr}
\toprule
side & minimum $|r_i|$ & median $|r_i|$ & maximum $|r_i|$\\
\midrule
left & $0.02610$ & $0.04562$ & $0.06514$\\
right & $0.01107$ & $0.04532$ & $0.07956$\\
\bottomrule
\end{tabular}
\end{center}
All residues are well separated from floating underflow and from the
$10^{-12}$ visibility threshold.  Thus no selected pole is canceled in the
scalar source-observation transfer function at this anchor.

Residue magnitude alone is not a coordinate-free angle.  Define the
normalized mode overlap
\begin{equation}\label{eq:mode-overlap}
 \eta_i=\frac{|r_i|}{\norm{X_i}_F\norm{Y_i}_F}.
\end{equation}
The minimum is $5.458\times10^{-3}$ on the left and
$1.125\times10^{-3}$ on the right.  Its reciprocal is the optimal norm
product for that single labeled mode pair.  The complete four-space
condition also reflects nonorthogonality among modes on each side.

\section{Canonical four-space condition}

The cross singular values are:
\begin{center}
\small
\begin{tabular}{lrrrr}
\toprule
side & $\sigma_1$ & $\sigma_2$ & $\sigma_3$ & $\sigma_4$\\
\midrule
left & $0.006738$ & $0.006565$ & $0.005549$ & $0.005207$\\
right & $0.002516$ & $0.002104$ & $0.001158$ & $0.001052$\\
\bottomrule
\end{tabular}
\end{center}

\begin{proposition}[Finite physical transversality]\label{prop:transverse}
Both canonical quartet pairings are nonsingular.  Their optimal
biorthogonal norm products and balanced frame norms are
\begin{center}
\begin{tabular}{lrr}
\toprule
side & $\chi_{\rm can}=1/\gamma$ & $\sqrt{\chi_{\rm can}}$\\
\midrule
left & $192.0508$ & $13.8582$\\
right & $950.2581$ & $30.8263$\\
\bottomrule
\end{tabular}
\end{center}
\end{proposition}

This proves finite transversality and rejects a stronger informal hope that
the exact spectral packet would be nearly orthogonal.  The right channel is
intrinsically close to a tangential source-observation pairing.

\section{Nonnormality of the base modes}

For normalized left/right eigenvectors $w_i^*v_i=1$, the base spectral
projector norm is
\begin{equation}\label{eq:projector-norm}
 \norm{P_i}=\norm{v_i}\norm{w_i}.
\end{equation}
The selected values range up to $17.54$ on the left and $17.69$ on the
right.  Thus the edge quartet is moderately nonnormal even before the source
and observation are applied.

The much larger packet condition, especially on the right, comes from the
combined geometry of spectral nonnormality, source activation, observation
activation, and near-tangency of the two four-spaces.  It cannot be inferred
from eigenvalue spacing alone.

\section{Comparison with temporal conditioning}

Let $\chi_t$ be the RH-185 oblique condition number of an accepted temporal
window.  Compare it with the exact spectral optimum through
\begin{equation}\label{eq:ratio}
 R_t=\frac{\chi_t}{\chi_{\rm can}}.
\end{equation}
The spaces differ, so $R_t$ can lie below one; the canonical optimum is an
optimum only for the canonical pair of spaces.

At the latest windows the values are approximately
\begin{center}
\begin{tabular}{lrrr}
\toprule
side & start & $\chi_t$ & $R_t$\\
\midrule
left & 18 & $193.974$ & $1.0100$\\
right & 18 & $956.830$ & $1.0069$\\
\bottomrule
\end{tabular}
\end{center}
Both relative discrepancies are below two percent.

\begin{proposition}[Finite geometry convergence signal]\label{prop:geometry-convergence}
At the latest accepted windows, temporal subspace gaps are below $0.051$ and
temporal oblique conditions differ from the canonical spectral optima by
less than two percent on both sides.
\end{proposition}

The proposition is a finite conjunction, not an asymptotic theorem.  It
shows that two independent diagnostics---subspace angle and cross-angle
condition---point to the same spectral endpoint.

\section{Why clipping is not the response}

RH-187 studies singular-value clipping of the temporal cross Gram and proves
the exact defect/conditioning tradeoff \citep{WangRH187}.  Clipping the small
singular values improves coordinate norms only by sacrificing exact
biorthogonality, and no audited clipping level closes the old coarse gate.

The present audit explains why.  Small cross singular values persist in the
exact canonical spectral spaces.  They are not solely noisy temporal
directions that should be deleted.  A later successful route should instead
use quantities invariant under balanced gauge, such as eigenvalues,
determinants, traces, and self-energy products in correctly reduced spaces.

\section{What would be needed uniformly}

A transportable all-level packet would need estimates of the form
\begin{equation}\label{eq:uniform-needs}
 \inf_k\min_i|r_{k,i}|>0,
 \qquad
 \inf_k\sigma_{\min}(Q_{L,k}^*Q_{R,k})>0,
\end{equation}
or a renormalized replacement in which their decay is explicitly canceled.
The present single-anchor numbers do not support either conclusion.

Indeed, a condition near $950$ leaves little room for untracked operator
perturbations.  Interval validation at one level is feasible, but a direct
uniform perturbation argument may be too expensive without additional
structure.

\section{Next use of the canonical packet}

RH-198 will audit the window-by-window approach to the canonical spaces and
fit only descriptive finite decay diagnostics.  RH-199 will then compare the
temporal determinant and power traces with their exact four-mode values.
These quantities may converge despite large frame norms because they are
similarity invariant.

The route remains inside Gate A.  Finite nonzero residues do not construct a
self-adjoint operator, prove a $T\log T$ law, recover prime-power weights, or
identify a zeta divisor.

\bibliographystyle{plainnat}
\bibliography{references}
\end{document}
